% Generated draft: requires mathematical and source review.
\documentclass[11pt]{article}
\usepackage[margin=25mm]{geometry}
\usepackage{fontspec}
\setmainfont{Latin Modern Roman}
\usepackage{amsmath,amssymb,mathtools,unicode-math}
\setmathfont{Latin Modern Math}
\usepackage{enumitem,needspace,xurl,hyperref}
\setlength{\parindent}{0pt}
\setlength{\parskip}{5pt}
\setlength{\emergencystretch}{3em}
\newcommand{\N}{\mathbb{N}}
\newcommand{\Z}{\mathbb{Z}}
\newcommand{\Q}{\mathbb{Q}}
\newcommand{\R}{\mathbb{R}}
\newcommand{\C}{\mathbb{C}}
\newcommand{\F}{\mathbb{F}}
\newcommand{\A}{\mathbb{A}}
\newcommand{\PP}{\mathbb{P}}
\newcommand{\E}{\mathbb{E}}
\newcommand{\eps}{\varepsilon}
\newcommand{\dd}{\,\mathrm{d}}
\begin{document}
\section*{UnsolvedMath ID: 2610}
\textbf{Source problem code:} catalogue entry
\subsection{Definitions and mathematical statement}
A finite group $G$ is almost simple if there is a nonabelian finite simple group $S$ with $S\leq G\leq\operatorname{Aut}(S)$, identifying $S$ with its inner automorphism group. The simple group $S$ is the socle, while $G$ may include outer automorphisms. An involution is an element of order exactly two.

The source presents two classification questions with the phrase ``in other words.'' They require a distinction.

\textbf{Involutory generation target.} Classify all finite almost simple groups $G$ for which there are involutions $\rho_0,\rho_1,\rho_2\in G$ satisfying
\[
 G=\langle\rho_0,\rho_1,\rho_2\rangle,\qquad
 \rho_0\rho_2=\rho_2\rho_0.
\]
The commuting pair can be named $\rho_0,\rho_2$ by relabeling. This is the displayed three-involution generation question, often described as $(2\times2,2)$-generation when the commuting pair generates a Klein four group.

\textbf{Regular rank-three polytope target.} Under the standard abstract-polytope interpretation, a regular polytope has an automorphism group transitive on its flags (maximal chains of incident faces). A rank-three polytope has vertices, edges and two-dimensional faces, together with its least and greatest faces; regularity is not a demand for a particular Euclidean realization. The standard algebraic target is a rank-three string C-group representation. Besides the generation and commuting relations above, it requires the intersection property
\[
 \langle\rho_i:i\in I\rangle\cap\langle\rho_i:i\in J\rangle
 =\langle\rho_i:i\in I\cap J\rangle
 \quad\text{for every }I,J\subseteq\{0,1,2\},
\]
with the empty generated subgroup equal to $\{1\}$. Classify the finite almost simple groups admitting such a triple.

The short generation formula omits this intersection condition. A particular generating triple does not acquire it merely from commutation. Whether the two existential classifications coincide in the intended class needs an additional argument, not provided by ``in other words.'' The source also does not explicitly distinguish abstract from geometric polytopes or state any additional nondegeneracy conventions. The two precise targets are therefore kept separate pending review.

The background says the finite simple case has been answered, referring to Archive Problem 7.30, and lists Ya. N. Nuzhin in its discussion/literature field. The literature triage dated 2026-08-17 reports a classification modulo the classification of finite simple groups, recording the sporadic exceptions $M_{11},M_{22},M_{23}$ and $\mathrm{McL}$ together with alternating and Lie-type results. It does not reproduce a full classification table or resolve which target each entry addresses; no blanket list is reconstructed.

The review also reports exact three-involution generation criteria for $\operatorname{GL}_n(q)$ and $\operatorname{PGL}_n(q)$ by Markovskaya--Nuzhin. General linear groups need not themselves be almost simple, so these are contextual family results, not a classification of all target groups. It cites Nuzhin's 2026 conditions linking triples to regular rank-three polytopes and Leemans's survey of string C-group representations. Exact criteria and extra hypotheses are not reproduced. The review says all intermediate extensions $S\leq G\leq\operatorname{Aut}(S)$ have not been completely classified in the sources checked.

The registry label is open and the triage label partially solved, with the distinction between generator and polytope targets explicitly unresolved in the supplied review. No source, PDF, classification table or proof was independently consulted, and no present-day status assessment is claimed.
\subsection{Short English statement}
Classify finite almost simple groups generated by three involutions with a commuting pair, and classify those that act as full automorphism groups of regular rank-three abstract polytopes. The latter standard formulation also needs the string C-group intersection property, absent from the source's claimed equivalence. Known simple-group and matrix-family results do not supply a complete almost-simple classification here.
\subsection{Sources}
\begin{itemize}
\item UnsolvedMath local registry, data/unsolvedmath/problems.json, permanent ID 2610, source code catalogue entry, ``Kourovka Notebook Problem 21.101,'' including both formulations and literature triage dated 2026-08-17; the discussion/literature field names Ya. N. Nuzhin.
\item \emph{Kourovka Notebook}, New Problems, 21st issue (2026), Problem 21.101, and Archive Problem 7.30, supplied link \url{https://alglog.org/21tkt.pdf}. The review cites the almost-simple question and known finite-simple case; the PDF was not opened here.
\item I. A. Markovskaya and Ya. N. Nuzhin, ``On generation of the groups $\operatorname{GL}_n(q)$ and $\operatorname{PGL}_n(q)$ by three involutions, two of which commute,'' \emph{Trudy Instituta Matematiki i Mekhaniki} 31 (2025), \url{https://doi.org/10.21538/0134-4889-2025-31-4-fon-03}. The review cites necessary and sufficient family criteria.
\item Ya. N. Nuzhin, ``Generating Sets of Involutions of Almost Simple Groups and Their Applications,'' \emph{Algebra and Logic} (2026), \url{https://doi.org/10.1007/s10469-026-09809-5}. The review cites additional conditions for polytope applications and further open problems.
\item Dimitri Leemans, ``String C-group representations of almost simple groups: a survey,'' \emph{Contemporary Mathematics} 764 (2021), 157--178, \url{https://doi.org/10.1090/CONM/764/15335}.
\item Ulam AI and the UnsolvedMath Contributors, \emph{UnsolvedMath: A Curated Collection of Open Mathematics Problems}. Dataset curation and original metadata are supplied under CC BY 4.0; referenced source material retains its own terms. No external source was independently consulted for this draft.
\end{itemize}
\end{document}
